% =========================================================
% Nederpelt & Geuvers, Type Theory and Formal Proof (2014)
% Sections 1.1--1.2 -- Input-output behaviour; the essence of functions
% Notes by William Sollers, typeset in the AIM node format.
%
% Page numbers: printed page = PDF page - 29.
% Fixes from corrections.md (items 1--7) are applied; the
% handwritten originals are unchanged.
% Slots marked \pending{...} are yours to write.
% =========================================================
\documentclass[11pt]{article}
\usepackage{../aim-notes}

\newcommand{\bred}{\to_\beta}

\title{Nederpelt \& Geuvers \S\S1.1--1.2: Functions, informally}
\author{William Sollers}
\date{Notes, draft 1}

\begin{document}
\maketitle

\begin{remark}[Status of these sections]
Everything here is \emph{informal}. Sections~1.1--1.2 motivate the
formal definitions that follow: $\Lambda$ (Def.~1.3.2), substitution
(\S1.6), and $\beta$-reduction (Def.~1.8.1). Each item below points
forward to its formal version.
\end{remark}

% ---------------------------------------------------------
\section{Input-output behaviour of functions (\S1.1)}

\begin{remark}[Functions compute]
\source{N\&G \S1.1, p.~1 (PDF p.~30), \textbf{verify} page.}
\begin{itemize}
  \item Input values are used to calculate output values.
  \item Variables in the input act as arbitrary (abstract) values.
  \item Concrete values replace the abstract values in order to compute.
\end{itemize}
\end{remark}

\begin{convention}[$\lambda$-notation]\label{conv:lambda-notation}
\source{N\&G \S1.1, p.~1.}
$\lambda x.\, x^2+1$ means $x \mapsto x^2+1$: ``the function mapping $x$
to $x^2+1$''.
\end{convention}

\begin{remark}[What the $\lambda$-calculus is]
\source{N\&G \S1.1, p.~1.}
The \emph{$\lambda$-calculus} is a formalisation of the basic aspects of
functions, in particular their construction and use. (This is a
description, not a definition; correction~1.)
\end{remark}

% ---------------------------------------------------------
\section{The essence of functions (\S1.2)}

\begin{remark}[Two construction principles, one evaluation rule]
\source{N\&G \S1.2, \textbf{verify} page.}
Construction: abstraction and application. Evaluation: $\beta$-reduction.
\end{remark}

\begin{definition}[Abstraction, informal]\label{def:abstraction-informal}
\source{N\&G \S1.2.}
From an expression $M$ and a variable $x$ we construct a new expression
$\lambda x.\, M$, called the \emph{abstraction of $x$ over $M$}.
Formal version: Def.~1.3.2(3).
\end{definition}

\begin{definition}[Application, informal]\label{def:application-informal}
\source{N\&G \S1.2.}
From expressions $M$ and $N$ we construct the expression $MN$, called the
\emph{application of $M$ to $N$}. Parentheses are added where necessary.
Formal version: Def.~1.3.2(2).
\end{definition}

\begin{example}
\begin{itemize}
  \item Application of $\lambda x.x$ to $\lambda y.y$ yields
        $(\lambda x.x)(\lambda y.y)$.
  \item Application of $f$ to $c$ yields $fc$, also written $f(c)$.
\end{itemize}
\end{example}

\begin{remark}[Meaningless expressions]
The construction principles also allow expressions with no evident meaning,
such as $xx$ (self-application). Ignored for now.
\end{remark}

\begin{remark}[Understanding: a function versus its output]
``Square'' is $\lambda x.\, x^2$ over $\mathbb{N}$, the map
$x \in \mathbb{N} \mapsto x^2 \in \mathbb{N}$. ``Square'' is not $x^2$:
$x^2$ is the abstract output value; ``square'' is the name of the function.
\end{remark}

\begin{remark}[Piecewise functions]
$\lambda$-notation is awkward for functions such as
$x \mapsto \begin{cases} x & x \ge 0\\ -x & x<0 \end{cases}$.
\textbf{Verify} whether this is N\&G's remark or your own (correction~4).
\end{remark}

% ---------------------------------------------------------
\subsection{Evaluation: $\beta$-reduction}

\begin{definition}[Substitution, informal]\label{def:subst-informal}
\source{N\&G \S1.2.}
$M[x := N]$ is $M$ with $N$ substituted for the \emph{free} occurrences of
$x$ (correction~2). Formal version: \S1.6.
\end{definition}

\begin{definition}[$\beta$-reduction, informal]\label{def:beta-informal}
\uses{def:subst-informal}
\source{N\&G \S1.2.}
$\beta$-reduction is the formalisation of the process of evaluating a
function. An expression of the form $(\lambda x.\, M)N$ can be rewritten to
$M[x := N]$; this is the \emph{$\beta$-reduction} of $(\lambda x.\, M)N$ to
$M[x:=N]$, written $(\lambda x.\, M)N \bred M[x:=N]$.
Formal version: Def.~1.8.1.
\end{definition}

\begin{example}\label{ex:beta-informal}
\uses{def:beta-informal}
One $\beta$-step, then evaluate the substitution (correction~3):
\begin{enumerate}[label=(\roman*)]
  \item $(\lambda y.\,5)(3) \bred 5[y:=3] \equiv 5$: replace $y$ by $3$ in
        the expression ``$5$''.
  \item $(\lambda x.\,x)(\lambda y.\,y) \bred x[x:=\lambda y.y] \equiv
        \lambda y.\,y$: replace $x$ by $\lambda y.y$ in ``$x$''.
  \item $(\lambda x.\,x^2+1)(3) \bred (x^2+1)[x:=3] \equiv 3^2+1$: replace
        $x$ by $3$ in ``$x^2+1$''.
\end{enumerate}
\end{example}

\begin{remark}[Compatibility, informal]\label{rem:compatibility-informal}
\uses{def:beta-informal}
\source{N\&G \S1.2, Example~1.2.4, p.~4 (\textbf{verify} number).}
Reduction is also possible on suitable \emph{parts} of expressions: when
$(\lambda x.\, M)N$ is a subexpression, it may be rewritten to $M[x:=N]$,
provided the rest of the expression is left unchanged. The rules describing
how reduction extends from subexpressions to bigger ones are the
\emph{compatibility rules}. Formal version: Def.~1.8.1.
\end{remark}

\begin{quantified}
Informally, for all expressions $M, N, L$ and every variable $x$:
\[
M \bred N \;\Longrightarrow\;
ML \bred NL,\quad LM \bred LN,\quad \lambda x.M \bred \lambda x.N .
\]
\pending{Compare with Def.~1.8.1 when you get there: which clauses does it
have, and is anything added or dropped?}
\end{quantified}

\begin{example}
By compatibility, $\lambda z.\,((\lambda x.x)(\lambda y.y)) \bred
\lambda z.\,(\lambda y.y)$. The result is the map $z \mapsto (y \mapsto y)$:
a function taking any input $z$ and returning the function $y \mapsto y$.
\end{example}

% ---------------------------------------------------------
\subsection{Functions of several variables}

\begin{remark}[Currying]\label{rem:currying}
\source{N\&G \S1.2, \textbf{verify} page.}
$\lambda$ binds exactly one variable. A function of two variables,
$f(x,y)$, takes an ordered pair; $\lambda$-notation instead takes the
arguments one after another. The two-argument function $(x,y) \mapsto M$
becomes
\[
  \lambda x.\, (\lambda y.\, M),
\]
called its \emph{Curried version} (correction~7). The form
$\lambda(x,y).M$ is not a $\lambda$-term (correction~6).
\end{remark}

\begin{remark}[Lean]
Lean's \texttt{fun x y => M} is notation for
\texttt{fun x => fun y => M}: Lean functions are Curried in exactly this
sense.
\end{remark}

\end{document}
